\begin{sol}[Esfera rodando dentro de un cuenco cilíndrico]

\subsubsection*{Restricción de rodadura}

El centro $G$ de la esfera se mueve sobre un círculo de radio $R-r$ centrado en $O$, con rapidez $v_G = (R-r)\dot\phi$. Usando la base polar $\{\hat{e}_r,\hat{e}_\theta\}$ centrada en $O$, la condición de rodadura sin deslizar en el punto de contacto $C$ (velocidad nula del punto material de la esfera en contacto) da:

$$v_C = (R-r)\dot\phi + \omega r = 0 \quad\Rightarrow\quad |\omega| = \frac{R-r}{r}\dot\phi$$

donde $\omega$ es la velocidad angular de \emph{spin} de la esfera (vista desde el marco fijo).

\subsubsection*{Pregunta 1: Ecuación de movimiento por energía}

La energía cinética combina traslación del centro de masa y rotación respecto a $G$, con $I_G = \frac{2}{5}mr^2$ para una esfera sólida:

$$T = \frac{1}{2}mv_G^2 + \frac{1}{2}I_G\omega^2 = \frac{1}{2}m(R-r)^2\dot\phi^2 + \frac{1}{2}\cdot\frac{2}{5}mr^2\cdot\left(\frac{R-r}{r}\right)^2\dot\phi^2$$

$$T = \frac{7}{10}m(R-r)^2\dot\phi^2$$

La energía potencial gravitatoria, con referencia en $O$:
$$V(\phi) = -mg(R-r)\cos\phi$$

Como el sistema es conservativo, $T + V = \text{const.}$ Derivando respecto al tiempo y dividiendo por $\dot\phi$ (asumiendo $\dot\phi \neq 0$):

$$\frac{7}{5}m(R-r)^2\ddot\phi + mg(R-r)\sin\phi = 0$$

$$\boxed{\ddot\phi + \frac{5g}{7(R-r)}\sin\phi = 0}$$

Para oscilaciones pequeñas, $\sin\phi \approx \phi$:
$$\ddot\phi + \omega_0^2\phi = 0, \qquad \omega_0^2 = \frac{5g}{7(R-r)}$$

\subsubsection*{Pregunta 2: Período de oscilación}

$$\boxed{T_{\text{esfera}} = 2\pi\sqrt{\frac{7(R-r)}{5g}}}$$

El período de un péndulo simple de longitud $\ell = R-r$ es $T_{\text{simple}} = 2\pi\sqrt{\ell/g}$. La razón entre ambos períodos es:

$$\frac{T_{\text{esfera}}}{T_{\text{simple}}} = \sqrt{\frac{7}{5}} \approx 1.18$$

\subsubsection*{Pregunta 3: Interpretación física}

El período de la esfera es aproximadamente un $18\%$ mayor que el del péndulo simple equivalente. Esto ocurre porque, a diferencia de una masa puntual que solo tiene energía cinética de traslación, la esfera rodante debe repartir la energía potencial liberada entre energía cinética de traslación \emph{y} de rotación (factor $2/5$ adicional por el momento de inercia). Como una fracción de la energía disponible se destina a hacer girar la esfera en torno a su propio eje, el centro de masa acelera más lentamente que en el caso de una masa puntual, resultando en oscilaciones más lentas y, por lo tanto, un período mayor.

\end{sol}
